% Source-level scope lemma by /root/singular_fibres.
% GPT-6 Astra Pro; reasoning effort not exposed. 2026-11-21.
\subsection{SF-Q6-03: the asymptotic size constructor does not specialize to the fixed two-step ranks}

\paragraph{Statement SF-Q6-03.}
For the rank data
\[
  (k,m,t,r,f)=(8,2,15,9,8),
\]
there is no scale $s$ satisfying \texttt{GlobalTheory.EX.RankVal s R}.
More generally, any rank data with $m=2$ satisfying
\texttt{RankVal s R} must have $s.x<81$.

\begin{proof}
The last clause of \texttt{RankVal}, defined in
\texttt{Development/AbelianVarietyTheory/Definitions/GlobalTwisting/
PairRowIffQdRows.lean}, is the equality of nonnegative integers
\[
 r+t=8fr+4+\texttt{TOv}(s)\,\texttt{nWv}(R,s).
\]
For the displayed ranks the left side is $24$ and the right side is
at least $8\cdot8\cdot9+4=580$, a contradiction.

For the second assertion, \texttt{RankVal} also requires
\[
 \sqrt{x\delta_1}\le m,\qquad
 \delta_1=x^{-\varepsilon/8},\qquad
 x>1,\quad 0<\varepsilon<1/52.
\]
Since $1-\varepsilon/8>1/2$, one has
$x\delta_1=x^{1-\varepsilon/8}\ge\sqrt{x}$.
Consequently $m\ge x^{1/4}$.  If $x\ge81$ this gives $m\ge3$,
contrary to $m=2$.  The archive proves precisely this general
implication in
\texttt{Proofs/IwasawaDecomposition/EventuallyNodeGapGeRankPow.lean},
\texttt{NodeGap.le\_m}: $M^4\le x\Rightarrow M\le m$.
\end{proof}

\paragraph{Scope for a global specialization.}
The actual constructor
\texttt{Proofs/GlobalInterpolation/ExistsExplicitParameters.lean},
\texttt{exists\_sizeData\_explicit}, assumes \texttt{RankVal};
the assembled results in
\texttt{Proofs/IwasawaDecomposition/SizeParametersNonempty.lean}
reach it through \texttt{Ctx.ok.rank}.  These theorems therefore
cannot fill the missing size datum for the fixed ranks of SF-Q6-02.
An eventual-in-$x$ theorem with this rank window cannot retain
$m=2$ along an unbounded scale family either.

Neither conclusion asserts that the abstract \texttt{SizeData$_2$}
or \texttt{DivInputs} types forbid every $m=2$ specialization.
Those types do not contain \texttt{RankVal} as a universal axiom.
A direct construction with different weights and independently
verified bounds would be a different route.

\paragraph{Exact locators and checks.}
The defining rank equality is in
\texttt{PairRowIffQdRows.lean}, lines 17--32;
\texttt{ChOK.rank} is in
\texttt{Definitions/GlobalTwisting/SetUnipotentTranslate.lean},
lines 36--39; and \texttt{NodeGap.le\_m} is in
\texttt{EventuallyNodeGapGeRankPow.lean}, lines 20--35.
The author read these complete definition/proof passages and the
full \texttt{exists\_sizeData\_explicit} constructor.  These are
human source-level proofs; no Lean compilation is claimed.
